\ExerciseSection

\begin{enumerate}
\def\labelenumi{\arabic{enumi})}
\tightlist
\item
  A positive real-valued function \(f\) on \(\mathbf{X} \times \mathbf{X}\) is a pseudometric on \(\mathbf{X}\) if and only if \(f(x, x) = 0\) for all \(x \in \mathbf{X}\) and
\end{enumerate}

\[
f (x, y) \leqslant f (x, z) + f (y, z)
\]

for all \(x, y, z\) in \(\mathbf{X}\) .

\begin{enumerate}
\def\labelenumi{\arabic{enumi})}
\setcounter{enumi}{1}
\tightlist
\item
  Let \(f\) be a mapping of \(\mathbf{X} \times \mathbf{X}\) into \([0, +\infty]\) . Then the family of sets \(\overline{f}^{1}([0, a])\) , where \(a\) runs through the set of all real numbers \(>0\) , forms a fundamental system of entourages of a uniformity on \(\mathbf{X}\) if and only if \(f\) satisfies the following conditions: \(a) f(x, x) = 0\) for all \(x \in \mathbf{X}\) ; \(b)\) for each \(a > 0\) there exists \(b > 0\) such that the relation \(f(x, y) \leqslant b\) implies \(f(y, x) \leqslant a\) ; \(c)\) for each \(a > 0\) there exists \(c > 0\) such that the relations \(f(x, z) \leqslant c\) and \(f(z, y) \leqslant c\) imply \(f(x, y) \leqslant a\) .
\end{enumerate}

The conditions \(b)\) and \(c)\) are satisfied in particular if there exists a mapping \(\varphi\) of \(I = [0, +\infty]\) onto itself which is continuous and zero at the point \(0\) , and a mapping \(\psi\) of \(I \times I\) into \(I\) which is continuous and zero at the point \((0, 0)\) , such that we have identically

\[
f (y, x) \leqslant \varphi (f (x, y)) \quad \text { and } \quad f (x, y) \leqslant \psi (f (x, z), f (z, y)).
\]

\begin{enumerate}
\def\labelenumi{\arabic{enumi})}
\setcounter{enumi}{2}
\tightlist
\item
  Let X be a topological space, let C be the topology of X, let \((f_{t})\) be a saturated family (no. 2) of pseudometrics on X, and let U be the uniformity defined by the family \((f_{t})\) .
\end{enumerate}

\begin{enumerate}
\def\labelenumi{\alph{enumi})}
\item
  The topology induced by \(\mathcal{U}\) is coarser than \(\mathfrak{C}\) if and only if the \(f_{t}\) are continuous on \(\mathbf{X} \times \mathbf{X}\) (with respect to the product of the topology \(\mathfrak{C}\) by itself).
\item
  The topology induced by \(\mathcal{U}\) is finer than \(\mathcal{C}\) if and only if, for each \(x_0 \in \mathbf{X}\) and each neighbourhood \(\mathbf{V}\) of \(x_0\) (in the topology \(\mathcal{C}\) ) there is an index \(\iota\) and a real number \(a > 0\) such that \(f_{\iota}(x_0, x) \geqslant a\) for all \(x \in \mathbb{C}\mathbf{V}\) .
\end{enumerate}

\begin{enumerate}
\def\labelenumi{\arabic{enumi})}
\setcounter{enumi}{3}
\tightlist
\item
  Let X be a non-Hausdorff uniformizable space and let R be the relation `` \(y \in \overline{\{x\}}\) '' between two points \(x, y\) of X.
\end{enumerate}

\begin{enumerate}
\def\labelenumi{\alph{enumi})}
\item
  Show that R is an equivalence relation on X, that every continuous mapping of X into a Hausdorff space is compatible with the relation R (Set Theory, R, § 5, no. 7), and that the quotient space X/R is completely regular.
\item
  If \(\mathcal{U}\) is any uniformity compatible with the topology of X, then the Hausdorff uniformity associated with \(\mathcal{U}\) (Chapter II, § 3, no. 9) is defined on the quotient space X/R and is compatible with the topology of X/R. The topological space X/R is called the completely regular space associated with the uniformizable space X.
\end{enumerate}

\begin{enumerate}
\def\labelenumi{\arabic{enumi})}
\setcounter{enumi}{4}
\tightlist
\item
  Let X be a uniformizable space. Show that the family of all pseudometrics on X which are continuous on \(X \times X\) defines a uniformity on X compatible with the topology of X. This uniformity \(U_{0}\) is called the universal uniformity on X; it is the finest of all uniformities compatible with the topology of X. If Y is any uniform space and if \(f: X \to Y\) is a continuous mapping, then f is uniformly continuous with respect to the universal uniformity on X; this is not the case for any other uniformity on X compatible with its topology. If there exists a uniformity U on X, compatible with its topology and such that X, endowed with this uniformity, is a complete space, then X is also a complete space with respect to any uniformity \(U'\) which is finer than U and coarser than \(U_{0}\) . (Observe that the Cauchy filters are the same for \(U'\) as for \(U_{0}\) .)
\end{enumerate}

¶ 6) a) Let X be a completely regular space and let U be a uniformity compatible with the topology of X. Let U* be the coarsest uniformity on X which makes uniformly continuous all the mappings of X into {[}0, 1{]} which are uniformly continuous with respect to U. Show that the uniformity U* is Hausdorff and compatible with the topology of X, and that X is precompact with respect to U*.

\begin{enumerate}
\def\labelenumi{\alph{enumi})}
\setcounter{enumi}{1}
\tightlist
\item
  Let K be a compact space. Show that every mapping \(f: X \to K\) which is uniformly continuous with respect to U is also uniformly continuous with respect to \(U^{*}\) (note that the unique uniformity on K is the coarsest for which all the continuous mappings of K into the interval {[}0, 1{]} are uniformly continuous). Hence show that \(U^{*}\) is the finest of the uniformities on X which are coarser than U and with respect to which X is precompact.
\end{enumerate}

\begin{enumerate}
\def\labelenumi{\arabic{enumi})}
\setcounter{enumi}{6}
\tightlist
\item
  Let X be a completely regular space. If \(U_{0}\) denotes the universal uniformity on X (Exercise 5), then \(U_{0}^{*}\) (Exercise 6) is the coarsest uniformity on X for which all continuous mappings of X into {[}0, 1{]} are uniformly continuous. Let \(\beta X\) denote the compact space obtained by completing X with respect to the uniformity \(U_{0}^{*}\) . \(\beta X\) is called the Stone-Čech compactification of X.
\end{enumerate}

\begin{enumerate}
\def\labelenumi{\alph{enumi})}
\item
  Show that every continuous mapping of X into a compact space K can be extended to a continuous mapping of \(\beta X\) into K (cf.~Set Theory, Chapter IV, § 3).
\item
  Let \(f\) be a homeomorphism of \(\mathbf{X}\) onto a dense subset \(\mathbf{X}'\) of a compact space \(\mathbf{K}\) , and let \(\overline{f}\) be the continuous mapping \(\beta \mathbf{X} \to \mathbf{K}\) which extends \(f\) . Show that \(\overline{f}(\beta \mathbf{X} - \mathbf{X}) = \mathbf{K} - \mathbf{X}'\) . {[}Suppose that there is a point \(a \in \beta \mathbf{X} - \mathbf{X}\) such that \(\overline{f}(a) = f(b)\) , where \(b \in \mathbf{X}\) ; consider a continuous mapping \(g: \beta \mathbf{X} \to [0, 1]\) such that \(g(a) = 1\) and \(g(b) = 0\) , and let \(P\) denote the set of all \(x \in \beta X\) such that \(g(x) \geqslant 1/2\) . There is a continuous mapping \(h: K \to [0, 1]\) such that \(h(f(b)) = 0\) and \(h(f(x)) = 1\) for all \(x \in P\) ; show that the relation \(h(f(a)) = 0\) contradicts the fact that \(a\) is in the closure of \(P\) in \(\beta X\) .{]}
\end{enumerate}

\begin{enumerate}
\def\labelenumi{\arabic{enumi})}
\setcounter{enumi}{7}
\tightlist
\item
  Let X be a completely regular space and let \(\beta X\) be its Stone-Cech compactification.
\end{enumerate}

\begin{enumerate}
\def\labelenumi{\alph{enumi})}
\item
  A filter \(\mathfrak{F}\) on X is said to be completely regular if \(\mathfrak{F}\) has a base \(\mathfrak{B}\) consisting of open sets such that for each set \(A \in \mathfrak{B}\) there exists a set \(B \in \mathfrak{B}\) contained in A, and a continuous mapping \(f: X \to [0, 1]\) which is equal to 0 on B and 1 on \(\mathbb{C}A\) . A completely regular filter is said to be maximal if there exists no strictly finer completely regular filter. Show that, if \(\mathfrak{F}\) is any completely regular filter on E, there is a maximal completely regular filter finer than \(\mathfrak{F}\) (use Zorn's lemma).
\item
  A completely regular filter \(\mathfrak{F}\) is maximal if and only if, for each pair of open sets A, B of X such that \(B \subset A\) and such that there is a continuous mapping of X into \([0, 1]\) which is equal to 0 on B and equal to 1 on \(\complement A\) , we have either \(A \in \mathfrak{F}\) , or else \(A \notin \mathfrak{F}\) there is a set of \(\mathfrak{F}\) which does not meet B (if all the sets of \(\mathfrak{F}\) meet B, consider the filter generated by the sets of \(\mathfrak{F}\) and the sets \(\overline{f}([0, a[)\) , where \(a \in ]0, 1[\) ).
\item
  Show that every maximal completely regular filter \(\mathfrak{F}\) is a Cauchy filter on \(\mathbf{X}\) with respect to the uniformity induced by that of \(\beta\mathbf{X}\) {[}show that each continuous mapping \(f: \mathbf{X} \to [\mathrm{0}, \mathrm{1}]\) has a limit with respect to \(\mathfrak{F}\) , by arguing by reductio ad absurdum and using \(b)\) {]}.
\end{enumerate}

\(d)\) Two distinct maximal completely regular filters \(\mathfrak{F},\mathfrak{F}^{\prime}\) cannot converge to the same point of \(\beta\mathbf{X}\) {[}observe that by \(b)\) there exists an open set \(\mathbf{A}\in\mathfrak{F}\) and an open set \(\mathbf{A}^{\prime}\in\mathfrak{F}^{\prime}\) such that \(\mathbf{A}\cap\mathbf{A}^{\prime}=\varnothing\) ; then argue by contradiction, using axiom \((\mathrm{O}_{\mathbf{IV}})\) {]}.

\begin{enumerate}
\def\labelenumi{\alph{enumi})}
\setcounter{enumi}{4}
\tightlist
\item
  The trace on X of the neighbourhood filter of any point \(x_{0} \in \beta X\) is a maximal completely regular filter B (argue by contradiction, using the definitions of a completely regular filter and of the neighbourhood filter of a point in \(\beta X\)) .
\end{enumerate}

\begin{enumerate}
\def\labelenumi{\arabic{enumi})}
\setcounter{enumi}{8}
\tightlist
\item
  \begin{enumerate}
  \def\labelenumii{\alph{enumii})}
  \tightlist
  \item
    Let X be a topological space, let \(\mathcal{G}\) be its topology and let \((f_{\iota})_{\iota \in \mathbb{I}}\) be a family of mappings of X into K = {[}0, 1{]} which are continuous in the topology \(\mathcal{G}\) . Let \(\mathcal{G}_0\) be the coarsest topology on X for which all the \(f_{\iota}\) are continuous. Show that \(\mathcal{G}_0 = \mathcal{G}\) provided that the family of sets \(\overline{f}_{\iota}^{1}([0, a[)\) , where \(\iota \in I\) and \(a \in ]0, 1[,\) is a subbase of \(\mathcal{G}\) (Chapter I, § 2, no. 3). The space X is then uniformizable, and if it is Hausdorff it is homeomorphic to a subspace of the cube K \(^{\mathbf{I}}\) .
  \end{enumerate}
\end{enumerate}

\begin{enumerate}
\def\labelenumi{\alph{enumi})}
\setcounter{enumi}{1}
\tightlist
\item
  Suppose that the family \((f_{t})\) is the family of all continuous mappings of X into K (with respect to C). Show that, if Y is any compact space, every mapping of X into Y which is continuous with respect to C is continuous with respect to \(C_{0}\) (embed Y in a cube). The topology C is uniformizable if and only if \(C_{0} = C\) .
\end{enumerate}

\begin{enumerate}
\def\labelenumi{\arabic{enumi})}
\setcounter{enumi}{9}
\tightlist
\item
  Let X be a completely regular space, let K be a compact subset of X, and let V be a neighbourhood of K in X.
\end{enumerate}

\begin{enumerate}
\def\labelenumi{\alph{enumi})}
\item
  Show that there is a continuous mapping of X into {[}0, 1{]} which is equal to 1 on K and equal to 0 on \(\complement V\) {[}use \((\mathrm{O}_{\mathbf{IV}})\), covering K by a finite number of suitably chosen neighbourhoods{]}.
\item
  Let \(X'\) be the quotient space of X obtained by identifying all the points of K. Show that \(X'\) is completely regular.
\end{enumerate}

¶ 11) Let X be a completely regular space. Two closed sets A, B in X are said to be completely separated if there exists a continuous mapping f of X into {[}0, 1{]} such that \(f(x) = 0\) on A and \(f(x) = 1\) on B. a) Let \(\beta X\) be the Stone-Čech compactification of X. Show that two closed sets A and B in X are completely separated if and only if their closures in \(\beta X\) do not intersect {[}use Exercises 7 a) and 10 a){]}.

\begin{enumerate}
\def\labelenumi{\alph{enumi})}
\setcounter{enumi}{1}
\tightlist
\item
  Let \(h\) be a homeomorphism of \(\mathbf{X}\) onto a dense subset \(\mathbf{X}'\) of a compact space \(\mathbf{K}\) , and let \(\overline{h}\) be the continuous mapping \(\beta \mathbf{X} \to \mathbf{K}\) which extends \(h\) . Suppose that, for each pair of completely separated closed subsets \(\mathbf{A}, \mathbf{B}\) of \(\mathbf{X}\) , the closures of \(\overline{h}(\mathbf{A})\) and \(h(\mathbf{B})\) in \(\mathbf{K}\) do not intersect. Show that \(\overline{h}\) is a homeomorphism of \(\beta \mathbf{X}\) onto \(\mathbf{K}\) . (Show that \(\overline{h}\) is injective: if \(a, b\) are distinct points of \(\beta \mathbf{X}\) such that \(\overline{h}(a) = \overline{h}(b)\) , consider a continuous mapping \(f: \beta \mathbf{X} \to [\mathbf{0}, \mathbf{1}]\) such that \(f(a) = 0\) , \(f(b) = 1\) , and consider the sets \(X \cap \overline{f}[0, 1/3]\) and \(X \cap \overline{f}[2/3, 1]\) .)
\end{enumerate}

¶ 12) Let X be a completely regular space and let βX be its Stone-Cech compactification.

\begin{enumerate}
\def\labelenumi{\alph{enumi})}
\item
  Let \(f\) be a finite continuous real-valued function on \(\beta X\) , such that \(f(x) > 0\) for all \(x \in X\) , and such that the set \(A = \overline{f}^1(0) \subset \beta X - X\) is not empty. Then there is a sequence \((a_n)\) of points of \(X\) such that the sequence of numbers \(\lambda_n = f(a_n) > 0\) is strictly decreasing and such that \(\lim_{n \to \infty} \lambda_n = 0\) . For each integer \(n > 0\) , let \(I_n\) be an open interval with centre \(\lambda_n\) in \(R_+^*\) , such that the closed intervals \(\overline{I}_n\) are mutually disjoint, and let \(M_n = \overline{f}^1(I_n) \cap X\) . For each subset \(H\) of \(N\) , let \(M_H\) be the union of the \(M_n\) for which \(n \in H\) ; for each filter \(\mathfrak{F}\) on \(N\) which is finer than the Fréchet filter, let \(\mathfrak{F}'\) denote the filter on \(X\) which has the \(M_H(H \in \mathfrak{F})\) as a base. Show that \(\mathfrak{F}'\) is a completely regular filter (Exercise 8), and that if \(\mathfrak{F}_1, \mathfrak{F}_2\) are two distinct ultrafilters on \(N\) , each of which is finer than the Fréchet filter, then there is no filter on \(X\) which is finer than each of \(\mathfrak{F}_1', \mathfrak{F}_2'\) . Deduce that Card \((A) \geqslant 2^{2\text{Card}(N)}\) {[}use Exercise 8 d) and Chapter I, § 4, Exercise 5{]}.
\item
  Suppose that X is infinite and discrete. Show that the uniformity induced on X by that of \(\beta\mathbf{X}\) is the uniformity of finite partitions (Chapter II, § 2, no. 2 and § 4, Exercise 12) {[}use Exercise 11 b){]}. Hence show that Card \((\beta\mathbf{X}) = 2^{2^{\operatorname{Card}(\mathbb{N})}}\) (cf.~Chapter I, § 4, Exercise 5). Show that there is a continuous real-valued function \(f \geqslant 0\) on X such that \(\overline{f}^{1}(0)\) is not empty and is contained in \(\beta\mathbf{X} - \mathbf{X}\) .
\item
  Under the same hypotheses as in b), show that if A is an infinite closed subset of \(\beta X\) , then \(\operatorname{Card}(A) \geqslant 2^{2\operatorname{Card}(N)}\) . {[}Show first that there exists an infinite sequence \((a_{n})\) of points of A, and for each n a neighbourhood \(V_{n}\) of \(a_{n}\) in \(\beta X\) , the \(V_{n}\) being mutually disjoint; deduce that every bounded real-valued function f defined on the set D of the \(a_{n}\) can be extended to a continuous function on \(\beta X\) ; for this, consider a real-valued function g defined on X and equal to \(f(a_{n})\) at each point of \(X \cap V_{n}\) . Hence conclude that \(\overline{D} \subset A\) is homeomorphic to the Stone-Čech compactification of D (cf.~§ 4, Exercise 17).{]}
\end{enumerate}

\begin{enumerate}
\def\labelenumi{\arabic{enumi})}
\setcounter{enumi}{12}
\item
  Let G be a locally compact, non-compact topological group. Show that the uniformity induced on the product space \(H = G \times G\) by the uniformity of its Stone-Čech compactification \(\beta H\) is strictly finer than the uniformity induced by that of \((\beta G) \times (\beta G)\) . {[}Assuming the result false, show, by applying Exercise 7 a) to the mapping \((x, y) \to xy^{-1}\) of H into G, that G would be isomorphic to a subgroup of a compact group; now apply Chapter III, § 3, no. 3, Proposition 4, Corollary 1.{]}
\item
  If a topological space X is such that every point of X has a closed neighbourhood which is a uniformizable subspace of X, show that X is uniformizable.
\end{enumerate}

¶ 15) a) Let X be a locally compact space and let \(\Phi\) be the set of all continuous mappings \(g: \mathbf{X} \to [\mathfrak{0}, \mathfrak{1}]\) which are zero on the complement of some compact set (depending on \(g\) ). Let \(\mathcal{U}_1\) be the coarsest uniformity on X for which all the mappings \(g \in \Phi\) are uniformly continuous. Show that \(\mathcal{U}_1\) is compatible with the topology of X, that the completion \(\hat{\mathbf{X}}\) of X with respect to \(\mathcal{U}_1\) is compact, and that \(\mathbf{X} - \hat{\mathbf{X}}\) consists of at most one point (show that if X is not compact, the filter of complements of relatively compact subsets of X is a Cauchy filter with respect to \(\mathcal{U}_1\) ).

\begin{enumerate}
\def\labelenumi{\alph{enumi})}
\setcounter{enumi}{1}
\item
  Show that \(\mathfrak{U}_1\) is the coarsest uniformity compatible with the topology of \(\mathbf{X}\) .
\item
  Conversely, let X be a completely regular space such that the set of uniformities compatible with the topology of X has a coarsest element \(U_{1}\) . Show that X is locally compact. (Using Exercise 6, show that X is precompact with respect to \(U_{1}\) ; then show that the complement of X in its completion with respect to \(U_{1}\) cannot have more than one point.)
\item
  Suppose that X is locally compact and \(\sigma\) -compact, and therefore the union of an increasing sequence \((\mathbf{U}_n)\) of relatively compact open sets such that \(\overline{\mathbf{U}}_n \subset \mathbf{U}_{n+1}\) (Chapter I, § 9, no. 9, Proposition 15). Show that there exists a continuous real-valued function \(f\) on \(\mathbf{X}\) , such that \(f(x) \leqslant n\) for \(x \in \overline{\mathbf{U}}_n\) and \(f(x) \geqslant n\) for \(x \in \mathbb{C} \overline{\mathbf{U}}_n\) {[}cf.~Exercise 10 a){]}. Let \(\mathcal{U}_2\) be the coarsest uniformity on \(\mathbf{X}\) for which \(f\) and all the functions \(g \in \Phi\) are uniformly continuous. Show that \(\mathcal{U}_2\) is compatible with the topology of \(\mathbf{X}\) and that there is an entourage \(\mathbf{V}\) of \(\mathcal{U}_2\) such that \(\mathbf{V}(x)\) is relatively compact in \(\mathbf{X}\) for all \(x \in \mathbf{X}\) (cf.~Chapter II, § 4, Exercise 9).
\end{enumerate}

¶ 16) Let X be a complete Hausdorff uniform space, and let U be the uniformity of X.

\begin{enumerate}
\def\labelenumi{\alph{enumi})}
\tightlist
\item
  Let A be a subset of X which is the union of a sequence \((\mathbf{F}_{n})\) of closed sets, and let \(x \notin A\) . Show that there is a continuous finite real-valued function \(f \geqslant 0\) on X, such that \(f(x) = 0\) and \(f(y) > 0\) for all \(y \in A\) . (If \(g_{n}\) is a continuous mapping of X into {[}0, 1{]} such that \(g_{n}(x) = 0\) and \(g_{n}(y) = 1\) for all \(y \in F_{n}\) , consider the function
\end{enumerate}

\[
f (x) = \sum_ {n = 0} ^ {\infty} g _ {n} (x) / 2 ^ {n}.
\]

\begin{enumerate}
\def\labelenumi{\alph{enumi})}
\setcounter{enumi}{1}
\tightlist
\item
  Let \(h \geqslant 0\) be a continuous real-valued function on \(\mathbf{X}\) , and let \(\mathbf{V}\) be the open set \(\overline{h}^{1}(]0, +\infty[)\) . Let \(\mathcal{U}'\) be the uniformity on \(\mathbf{V}\) which is the least upper bound of the uniformity induced by U and the uniformity defined by the pseudometric
\end{enumerate}

\[
r (x, y) = \left| \frac {\mathrm{1}}{h (x)} - \frac {\mathrm{1}}{h (y)} \right|.
\]

Show that V is complete with respect to \(U'\) .

\begin{enumerate}
\def\labelenumi{\alph{enumi})}
\setcounter{enumi}{2}
\tightlist
\item
  Let B be a subset of X which is the intersection of a family of sets \((\mathbf{A}_{\lambda})\) , each of which is the union of a countable family of closed sets. Show that there exists a uniformity on B which is compatible with the topology induced on B by the topology of X, and such that B is complete with respect to this uniformity {[}use b){]}.
\end{enumerate}

\begin{enumerate}
\def\labelenumi{\arabic{enumi})}
\setcounter{enumi}{16}
\item
  Let X be a topological space such that each \(x \in X\) has a fundamental system of neighbourhoods which are both open and closed. Show that X is uniformizable.
\item
  Let X be a countable completely regular space. Show that, for each \(x \in X\) , the neighbourhoods of x which are both open and closed form a fundamental system of neighbourhoods of x (cf.~§ 6, no. 4).
\item
  Let X be a locally compact space. Show that every function \(f \geqslant 0\) which is lower semi-continuous on X is the upper envelope of a family of continuous functions \(\geqslant 0\) , each of which is zero on the complement of a compact set.
\item
  \begin{enumerate}
  \def\labelenumii{\alph{enumii})}
  \tightlist
  \item
    Let X be a completely regular space, and let a be a point of X. Then there exists a continuous function \(f \geqslant 0\) on X such that \(f(a) = 0\) and \(f(x) > 0\) whenever \(x \neq a\) , if and only if there is a sequence \((\mathbf{V}_n)\) of neighbourhoods of a in X such that \(\bigcap_{n} V_n = \{a\}\) {[}cf.~Exercise 16 a){]}.
  \end{enumerate}
\end{enumerate}

\begin{enumerate}
\def\labelenumi{\alph{enumi})}
\setcounter{enumi}{1}
\tightlist
\item
  Let X be a completely regular space, each point of which has a countable fundamental system of neighbourhoods. Show that, in the Stone-Čech compactification \(\beta\mathrm{X}\) , the set X is equal to the set of points which have a countable fundamental system of neighbourhoods in \(\beta\mathrm{X}\) {[}use a) and Exercise 12{]}. Let \(\mathbf{X}'\) be another completely regular space, each point of which has a countable fundamental system of neighbourhoods. Show that if the Stone-Čech compactifications \(\beta\mathrm{X}\) and \(\beta\mathrm{X}'\) are homeomorphic, then X and \(\mathbf{X}'\) are homeomorphic.
\end{enumerate}

¶ 21) a) Let X be a topological space. Show that the following two properties are equivalent: α) every finite continuous real-valued function on X is bounded; β) every bounded continuous real-valued function on X attains its bounds. X is said to be pseudo-compact if it has these properties. If X is pseudo-compact and if f is any continuous mapping of X into a topological space \(X'\) , then \(f(X)\) is pseudo-compact.

\begin{enumerate}
\def\labelenumi{\alph{enumi})}
\setcounter{enumi}{1}
\item
  Consider the following two properties of a topological space X: \(\gamma\) if \((\mathbf{U}_n)\) is any countable open covering of X, then X is the union of a finite number of the closures \(\overline{\mathbf{U}}_n\) ; \(\delta\) every countable filter base on X which is formed of open sets has at least one cluster point. Show that \(\gamma\) and \(\delta\) are equivalent and that they imply that X is pseudo-compact. In particular, every absolutely closed topological space (Chapter I, § 9, Exercise 19) is pseudo-compact. If X has properties \(\gamma\) and \(\delta\) , then so does every subspace of X which is the closure of an open set of X {[}cf.~§ 4, Exercise 26 b){]}.
\item
  Consider the following property of a topological space \(\mathbf{X}\) : \(\zeta)\) every locally finite open covering of \(\mathbf{X}\) is finite. Show that \((\zeta) \Longrightarrow \gamma)\) . {[}If \((\mathbf{U}_n)\) is an increasing sequence of open subsets of \(\mathbf{X}\) which cover \(\mathbf{X}\) and are such that \(\mathbf{U}_{n+1} \notin \overline{\mathbf{U}}_n\) , consider the covering formed by the sets \(\mathbf{U}_{n+1} \cap \mathbb{C}\overline{\mathbf{U}}_n\) and the complement of a sequence \((a_n)\) such that
\end{enumerate}

\[
a _ {n + 1} \in \mathrm{U} _ {n + 1} \cap \left[ \overline {{\mathrm{U}}} _ {n}. \right]
\]

If X is regular, show that \(\gamma\Longrightarrow\zeta\) . {[}Let \((\mathrm{U}_{\alpha})\) be an infinite, locally finite, open covering of X. Define by induction a sequence \((\alpha_{n})\) of indices, a sequence \((x_{n})\) of points of X, and for each \(x_{n}\) two open neighbourhoods \(V_{n}\) and \(W_{n}\) of \(x_{n}\) such that: (i) we have \(\overline{V}_{n}\subset W_{n}\) , \(W_{n}\subset U_{\alpha_{n}}\) , and \(W_{n}\) meets only a finite number of sets \(U_{\alpha}\) ; (ii) \(U_{\alpha_{n}}\) meets none of the \(W_{k}\) with indices k\textless n.~Now consider the covering formed by the \(W_{n}\) and the complement of the union of the \(\overline{V}_{n}\) .{]}

\begin{enumerate}
\def\labelenumi{\alph{enumi})}
\setcounter{enumi}{3}
\item
  In a completely regular space X the properties \(\alpha\) , \(\beta\) , \(\gamma\) , \(\delta\) ) and \(\zeta\) are all equivalent, and are equivalent to the following property: \(\theta\) X is precompact in any uniformity compatible with its topology. {[}To show that \(\alpha) \Longrightarrow \zeta\) ), note that if \((\mathbf{U}_n)\) is a countably infinite, locally finite, open covering of X and if \(a_n \in \mathbf{U}_n\) , then there is a continuous real-valued function \(f\) on X such that \(f(a_n) = n\) for each \(n\) {[}use \((\mathrm{O}_{\mathrm{IV}})]\) . To show that \(\theta) \Longrightarrow \alpha\) ), use Exercise 5 and observe that in a precompact space every uniformly continuous real-valued function is bounded. To show that \(\gamma) \Longrightarrow \theta\) ), note that if X is not precompact with respect to a uniformity \(\mathcal{U}\) , then there is a symmetric entourage V of \(\mathcal{U}\) and a sequence \((x_n)\) of points of X such that no two of the sets \(\mathrm{V}(x_n)\) intersect.{]}
\item
  If X is a completely regular pseudo-compact space which is complete with respect to some uniformity compatible with its topology, then X is compact.
\item
  On a completely regular pseudo-compact space X, the universal uniformity (Exercise 5) is induced by the uniformity of the Stone-Čech compactification of X, and is the unique uniformity, compatible with the topology of X, for which all the continuous mappings of X into {[}0, 1{]} are uniformly continuous.
\end{enumerate}

\begin{enumerate}
\def\labelenumi{\arabic{enumi})}
\setcounter{enumi}{21}
\tightlist
\item
  Let X be a Hausdorff uniform space, let Y be a closed subset of X, and let f be a bounded, uniformly continuous, real-valued function on Y. Show that f has an extension \(\overline{f}\) to X which is bounded and uniformly continuous on X. {[}We may assume that \(f(Y) \subset [0, 1]\) . For each dyadic number \(r = k/2^{n} \in [0, 1]\) , let \(A(r)\) be the set of all \(x \in Y\) such that \(f(x) \leqslant r\) , and let \(B(r) = A(r) \cup (X - Y)\) . Define by induction * (as in the proof of Theorem 1 of § 4, no. 1) * an open set \(U(r)\) for each dyadic number \(r \in [0, 1]\) such that, whenever r and \(r'\) are two dyadic numbers and \(r < r'\) , there is an entourage V of the uniformity for X for which
\end{enumerate}

\[
\mathbf {V} (\mathbf {A} (r)) \subset \mathbf {U} (r ^ {\prime}), \quad \mathbf {V} (\mathbf {U} (r)) \subset \mathbf {U} (r ^ {\prime}), \quad \mathbf {V} (\mathbf {U} (r)) \subset \mathbf {B} (r ^ {\prime}).
\]

Then define \(\overline{f}(x)\) to be the greatest lower bound of the dyadic numbers \(r\) such that \(x \in \mathrm{U}(r).]\)

